% Generated draft: requires mathematical and source review.
\documentclass[11pt]{article}
\usepackage[margin=25mm]{geometry}
\usepackage{fontspec}
\setmainfont{Latin Modern Roman}
\usepackage{amsmath,amssymb,mathtools,unicode-math}
\setmathfont{Latin Modern Math}
\usepackage{enumitem,needspace,xurl,hyperref}
\setlength{\parindent}{0pt}
\setlength{\parskip}{5pt}
\setlength{\emergencystretch}{3em}
\newcommand{\N}{\mathbb{N}}
\newcommand{\Z}{\mathbb{Z}}
\newcommand{\Q}{\mathbb{Q}}
\newcommand{\R}{\mathbb{R}}
\newcommand{\C}{\mathbb{C}}
\newcommand{\F}{\mathbb{F}}
\newcommand{\A}{\mathbb{A}}
\newcommand{\PP}{\mathbb{P}}
\newcommand{\E}{\mathbb{E}}
\newcommand{\eps}{\varepsilon}
\newcommand{\dd}{\,\mathrm{d}}
\begin{document}
\section*{UnsolvedMath ID: 2974}
\textbf{Source problem code:} catalogue entry
\subsection{Definitions and mathematical statement}
Let $(X,\omega)$ be a closed symplectic $4$-manifold. A Lefschetz pencil here has a nonempty finite base set $B\subset X$ and a map $\pi:X\setminus B\to S^2$, with the standard positive complex local models $[z_1:z_2]$ at base points and $z_1^2+z_2^2$ at isolated critical points; regular fibers are symplectic. The fiber genus is that of the closed completion of a regular fiber, equivalently of the fiber after blowing up $B$. A Lefschetz fibration has no base points and is a different object.

The first question asks whether every such $X$ has at least two inequivalent pencils of one common fiber genus $g$ for sufficiently large $g$. The triage's all-large-genus reading is
\[
\forall(X,\omega)\ \exists g_0\ \forall g\geq g_0,\qquad
\text{there exist two inequivalent genus-}g\text{ pencils on }X.
\]
The source does not clarify whether $g$ ranges over every integer or only genera attainable by pencils. The second question asks for infinitely many inequivalent pencils at fixed genus on the same $X$, with the same possible ambiguity about all large genera versus some sufficiently large genus. It is not enough to obtain arbitrarily large finite families at different genera or on different blow-ups.

The supplied equivalence description uses a diffeomorphism of the total $4$-manifold commuting with the pencil maps. For two pencils, a literal fixed-base diagram would require $h(B_1)=B_2$ and $\pi_2\circ h=\pi_1$. A commonly used more flexible diagram permits a base diffeomorphism $u:S^2\to S^2$ and $\pi_2\circ h=u\circ\pi_1$. The notes do not explicitly distinguish these conventions; the choice matters, in particular for critical-value positions. Nor do they require $h$ to be a symplectomorphism. The intended equivalence must be fixed rather than replaced without notice.

The background notes that increasing degree in Donaldson's construction gives arbitrarily high genera on $X$, but also increases the number of base points. Blowing them up gives fibrations on different manifolds $X\#|B|\overline{\C\mathrm P}^{,2}$, not the fixed-total-space conclusion. Outside rational and ruled surfaces, it reports $|B|\leq2g-2$. Hence an infinite family at fixed genus would contain infinitely many with the same base-point count, producing inequivalent same-genus fibrations after blow-up; the converse need not produce pencils.

The notes report Baykur's arbitrarily large finite families, for nonrational/nonruled manifolds possibly after blow-ups, with common genus, common base-point count, and matching singular-fiber types, and a slightly weaker rational/ruled result. These are results up to birational modification, not an infinite fixed-$X$, fixed-genus family. They also cite examples on particular manifolds by Park--Yun, Baykur--Hayano, Baykur--Hayano--Monden, and Hamada, and say available monodromy invariants in those works distinguish only finitely many classes.

The registry and supplied triage dated 2026-08-17 retain the universal and infinite versions as open. All construction and limitation claims are attributed supplied reports, not independently verified results here.
\subsection{Short English statement}
Does every closed symplectic $4$-manifold support two, or infinitely many, inequivalent Lefschetz pencils of the same sufficiently large genus? The total manifold must stay fixed. Known finite-family results after blow-ups do not give the requested infinite family, and the exact genus and equivalence conventions need clarification.
\subsection{Sources}
\begin{itemize}
\item K3: A New Problem List in Low-Dimensional Topology, Problem 4.98 (2026): \url{https://aimath.org/pastworkshops/kirbylistrep.pdf}. This supplied citation is retained without opening or generating a PDF.
\item [Bay16] R. İnanç Baykur, ``Inequivalent Lefschetz fibrations and surgery equivalence of symplectic 4-manifolds,'' Journal of Symplectic Geometry 14(3) (2016), 671--686, DOI 10.4310/JSG.2016.v14.n3.a2.
\item [Bay19] R. İnanç Baykur, ``Inequivalent Lefschetz fibrations on rational and ruled surfaces,'' Breadth in contemporary topology, Proceedings of Symposia in Pure Mathematics 102 (2019), 21--28, American Mathematical Society, DOI 10.1090/pspum/102/02.
\item [PY09] Jongil Park and Ki-Heon Yun, ``Nonisomorphic Lefschetz fibrations on knot surgery 4-manifolds,'' Mathematische Annalen 345(3) (2009), 581--597, DOI 10.1007/s00208-009-0366-0.
\item [PY17] Jongil Park and Ki-Heon Yun, ``Lefschetz fibrations on knot surgery 4-manifolds via Stallings twist,'' Michigan Mathematical Journal 66(3) (2017), 481--498, DOI 10.1307/mmj/1497513628.
\item [BH16b] R. İnanç Baykur and Kenta Hayano, ``Multisections of Lefschetz fibrations and topology of symplectic 4-manifolds,'' Geometry and Topology 20(4) (2016), 2335--2395, DOI 10.2140/gt.2016.20.2335.
\item [BHM23] R. İnanç Baykur, Kenta Hayano, and Naoyuki Monden, ``Unchaining surgery and topology of symplectic 4-manifolds,'' Mathematische Zeitschrift 303(3) (2023), Paper 77, 32 pages, DOI 10.1007/s00209-023-03204-x.
\item [Ham17] Noriyuki Hamada, ``Sections of the Matsumoto-Cadavid-Korkmaz Lefschetz fibration'' (2017), arXiv:1610.08458, supplied identifier.
\item Local registry data/unsolvedmath/problems.json, permanent catalogue entry, source code catalogue entry, ``Kirby Problem 4.98''; statement, complete supplied literature notes, and triage dated 2026-08-17. Dataset attribution: Ulam AI and the UnsolvedMath Contributors.
\end{itemize}
\end{document}
